\documentclass[11pt]{article}
\usepackage{xltxtra}
\usepackage{bookmark}
\usepackage{hyperref}
\usepackage{multicol}
\usepackage{xcolor}
\usepackage{calc}
\usepackage{graphicx}
\usepackage{tikz}
\usepackage{fontspec}
\usepackage{xeCJK}
\usepackage{relsize}
\usepackage{xspace}
\usepackage{fontawesome5}
\usepackage{titlesec}
\usepackage{enumitem}
\usepackage{siunitx}
\usepackage{amssymb}
\usepackage{tabularx}
\usepackage{multicol}
\usepackage{fancybox}
\usepackage{float}
\hypersetup{hidelinks}
\usepackage{url}
\urlstyle{tt}
\usetikzlibrary{calc}

%%%%%%%%%%%%%%%%%%%%
% 设置
%%%%%%%%%%%%%%%%%%%%

\protected\def\Cpp{{C\nolinebreak[4]\hspace{-.05em}\raisebox{.28ex}{\relsize{-1}++}}\xspace}	% 这是个更好看的 C++ 写法，可以使用 \Cpp 来代替出文本的 C++

\setlength{\parindent}{0pt}							% 取消全局段落缩进
\pagenumbering{gobble}								% 取消页码显示
\setlist[itemize]{leftmargin=*}		
\setlist[enumerate]{leftmargin=*}	
% 这里把表格的行间距调成1.2倍了
\renewcommand{\arraystretch}{1.2}
% 这里把正文的行间距调成1.2倍了
\linespread{1.2}

\titleformat{\section}					    % 将原标题前面的数字取消了
  {\LARGE\bfseries\raggedright} 		      % 字体改为LARGE，bold，左对齐
  {}{0em}                      			  % 可用于添加全局标题前缀
  {}                           			  % 可用于添加代码
  [{\color{secondary_color}\titlerule}]     % 标题下方加一条线
\titlespacing*{\section}{0cm}{*1.2}{*1.2}	% 标题左边留白，上方1.2倍，下方1.2倍

\titleformat{\subsection}				    % 将原二级标题前面的数字取消了
  {\large\bfseries\raggedright} 		      % 字体改为large，bold，左对齐
  {}{0em}                      			  % 可用于添加全局二级标题前缀
  {}                           			  % 可用于添加代码
  []
\titlespacing*{\subsection}{0cm}{*1.2}{*1.2}% 二级标题左边留白，上方1.2倍，下方1.2倍
% 页面大小与页边距，按需求调整

\usepackage[
	a4paper,
	left=1.2cm,
	right=1.2cm,
	top=1.5cm,
	bottom=1cm,
	nohead
]{geometry}

% 英文字体
\setmainfont[
    Path=fonts/,
    Extension=.otf,
    BoldFont=*-Bold,
]{NotoSerifSC}
% 中文字体
\setCJKmainfont[
    Path=fonts/,
    Extension=.otf,
    BoldFont=*-Bold,
]{NotoSerifSC}

% 自定义颜色（参考 https://github.com/seumxc/SEU-Logo）
% 东大绿
\definecolor{primary_color}{RGB}{63, 97, 21}
% 东大黄
\definecolor{secondary_color}{RGB}{255, 204, 0}

%%%%%%%%%%%%%%%%%%%%
% 文章内容
%%%%%%%%%%%%%%%%%%%%

% 学院
\newcommand{\school}{自动化学院 | School of Automation} 

% 联系方式
\newcommand{\contact}{
    % 根据个人喜好选择字号
    % \small              % 小
    % \footnotesize       % 更小
    \scriptsize             % 再小一号
    \textcolor{white}{
        % 邮箱
        \faEnvelope \quad \href{mailto:youremail@seu.edu.cn}{youremail@seu.edu.com}
        \hspace{4em}
        % 手机号
        \faPhone \quad  130-6666-0000
        % 别的联系方式，如微信、GitHub等
        % \hspace{4em}
        % \faGithub \quad \href{https://github.com/xxxx}{https://github.com/xxxx}
    }
}

\begin{document}

    %%%%%%%%%%%%%%%%%%%%
    % 页眉、页脚和背景（如果有多页简历，请把页眉页脚和背景复制粘贴到第二页的内容之前）
    %%%%%%%%%%%%%%%%%%%%

    % 页眉：校标组合+学院名
    \begin{tikzpicture}[remember picture, overlay]
        \node[anchor=north, inner sep=0pt](header) at (current page.north){
            \includegraphics[width=\paperwidth]{images/header-darkgreen.png}
        };
        \node[anchor=west](school_logo) at (header.west){
            \hspace{0.5cm}
            \includegraphics[width=0.2\textwidth]{images/banner-white.png}
        };
        \node[anchor=east](school_name) at(header.east){
            \textcolor{white}{\textbf{\school}}
            \hspace{0.5cm}
        };
    \end{tikzpicture}
    \vspace{-3.5em}

    % 页脚，联系方式
    \begin{tikzpicture}[remember picture, overlay]
        \node[anchor=south, inner sep=0pt](footer) at (current page.south){
            \includegraphics[width=\paperwidth]{images/footer-darkgreen.png}
        };
        % 联系方式
        \node[anchor=center] at(footer.center){\contact};
    \end{tikzpicture}

    % 背景
    \begin{tikzpicture}[remember picture, overlay]
        \node[opacity=0.05] at(current page.center){
            \includegraphics[width=0.7\paperwidth, keepaspectratio]{images/logo-yellow.png}
        };
    \end{tikzpicture}

    %%%%%%%%%%%%%%%%%%%%
    % 简历正文
    %%%%%%%%%%%%%%%%%%%%

    \begin{figure}[h]
        % 左半边，信息，比例占行宽78%
        % 合并了个人信息和教育背景部分，以便插入更大的照片
        \begin{minipage}[t]{0.78\textwidth}
            % 个人信息
            % \faGraduationCap这类\fa开头的都是font awesome里的logo，想换成其他logo的话，可以看一下附带的fontawsome.pdf，自行替换。
            \section{\makebox[\widthof{\faAddressCard}][c]{\color{primary_color}{\faAddressCard}}\quad 个人信息}
            \begin{tabularx}{\linewidth}{p{\widthof{出生日期:}}Xp{\widthof{政治面貌:}}X}
                姓\ \ \ \ \ \ \ \ 名: & 你的名字 & 
                性\ \ \ \ \ \ \ \ 别: & 你的性别 \\
                出生年月: & 你的出生年月 & 
                政治面貌: & 你的政治面貌 \\
                %% 想多加几行的话，就按上面的格式自行补充
                %% 想加粗的话\textbf{}
                %% 想多加几列的话，把\begin{tabularx}{\textwidth}{这里}的内容改一下，可以自己搜一下tabularx怎么用，也可以问gpt/文心一言/讯飞。
            \end{tabularx}
            
            \vspace{0.2em}
    
            % 教育背景
            \section{\makebox[\widthof{\faGraduationCap}][c]{\color{primary_color}{\faGraduationCap}}\quad 教育背景}
            
            {\large \textbf{东南大学}}，专科 \hfill {江苏，南京} \\
            {{你所在的学院}}，专业：你的专业 \hfill {2000年9月-2010年6月} \\
            {主修课程}：课程1、课程2、课程3、课程4\ 等。
            
            \vspace{0.5em}
            {\large \textbf{东南大学}}，本科 \hfill {江苏，南京} \\
            {{你所在的学院}}，专业：你的专业 \hfill {2010年9月-2020年6月} \\
            \textbf{主修课程}：课程1、课程2、课程3、课程4\ 等。
            
            \vspace{0.5em}
            {\large \textbf{东南大学}}，硕士 \hfill {江苏，南京} \\
            {{某些学院}}，导师：\href{导师的个人主页.site}{导师名字}\ 导师职称 \hfill {2020年9月-至今} \\
            \textbf{研究方向}：方向1、方向2、方向3、方向4\ 等。
        \end{minipage}
        \hfill
        % 右半边，照片，比例占行宽20%
        \begin{minipage}[t]{0.2\textwidth}
            \vspace{2em}
            \setlength{\fboxsep}{0pt}
            \doublebox{\includegraphics[width=\linewidth]{images/avatar.png}}
        \end{minipage}
    \end{figure}
    \vspace{-1em}

    % 科研成果
    \section{\makebox[\widthof{\faGraduationCap}][c]{\color{primary_color}{\faGraduationCap}}\quad 科研成果}

    % 科研著作（研究生）
    This is One of Your Paper Published in Conference A. \\
    \textbf{Mingzi Nide}, Daoshi Nide. \hfill 
    发表于 \textbf{Conference A}（CCF-A类会议） 

    \vspace{0.5em}
    This is Another Paper. \\
    \textbf{Mingzi Nide}, Shidi Nide, Daoshi Nide. \hfill 
    发表于 \textbf{Conference B} （CCF-A类会议）

    \vspace{0.5em}
    This is A Journal Paper.\\
    \textbf{Mingzi Nide}, Shixiong Nide, Daoshi Nide. \hfill 
    发表于 \textbf{Journal C} （SCI-1区）

    % 项目经历\科研经历\项目与教学（标题请根据需要修改）
    \section{\makebox[\widthof{\faChalkboardTeacher}][c]{\color{primary_color}{\faChalkboardTeacher}}\quad 项目与教学}
    % \vspace{0.5em}
    {\large{\textbf{项目名称}}} \hfill {横向/纵向项目-已完结/进行中}\\
    \textbf{你在项目中扮演的角色} \hfill 2020年9月至2021年9月\\
    项目简介。

    \vspace{0.5em}
    {\large{\textbf{某某主题讨论班}}}，主讲/参与 \hfill {2020年夏季}\\
    主要内容：内容1，内容2，内容3\ 等。
    
    \vspace{0.5em}
    {\large{\textbf{课程名称}}}，助教 \hfill {2021年夏季}\\
    主要内容：内容1，内容2，内容3\ 等。
    
    % 如果每行的内容不是很多，可以考虑使用 minipage，将内容分列展示
    \begin{minipage}[t]{0.6\textwidth}
        \section{\makebox[\widthof{\faWrench}][c]{\color{primary_color}{\faWrench}}\quad 技能特长}
        \begin{itemize}
            \item 熟练使用 Python、Golang、Rust 等编程语言。
            \item 熟练使用 Tensorflow、Pytorch 等深度学习框架。
            \item 熟悉 Windows 与 Linux 端开发。
        \end{itemize}
    \end{minipage}
    \hfill
    \begin{minipage}[t]{0.35\textwidth}
        \section{\makebox[\widthof{\faStar}][c]{\color{primary_color}{\faStar}}\quad 兴趣爱好}
        \begin{itemize}
            \item 我喜欢唱
            \item 跳
            \item Rap
            \item 篮球
        \end{itemize}
    \end{minipage}
    
    % \newpage
    % % 如有需要，可以添加额外的页面。不要忘记添加页眉页脚和背景相关的代码。

    % % 竞赛经历
    % \section{\makebox[\widthof{\faTrophy}][c]{\color{primary_color}{\faTrophy}}\quad 竞赛经历}
    % \begin{table}[h!]
    %     \begin{tabularx}{\textwidth}{Xp{\widthof{第零负责人}}p{\widthof{国家级-第100名}}p{\widthof{2030年13月}}}
    %         \textbf{比赛1} & 第一负责人 & 国家级-第10名 & 2023年4月 \\
    %         \textbf{比赛2} & 个人参赛 & 国家级-一等奖 & 2023年8月\\
    %         \textbf{比赛3} & 个人参赛 & 省级-一等奖 & 2022年12月\\
    %         % 同理，可以自己加
    %     \end{tabularx}
    % \end{table}

    % % 技能特长
    % \section{\makebox[\widthof{\faWrench}][c]{\color{primary_color}{\faWrench}}\quad 技能特长}
    % \begin{itemize}
    %     \item 熟练使用\Cpp 、Python、Matlab编程语言。
    %     \item 熟悉Windows与Linux端开发。
    %     \item 熟练使用Tensorflow，Pytorch等深度学习框架。
    %     \item 熟练掌握\Cpp 与Python环境下OpenCV与Qt应用的开发，且熟练使用Qt Creator软件。
    %     \item 熟练使用Altium Designer与LCEDA进行封装绘制与板子设计。
    %     \item 熟练使用Keil，Arduino IDE等集成开发软件。
    %     \item 了解模式识别，强化学习，遗传算法，知识蒸馏等相关概念。
    % \end{itemize}

    % % 所获荣誉
    % \section{\makebox[\widthof{\faStar}][c]{\color{primary_color}{\faStar}}\quad 所获荣誉}
    % \begin{multicols}{2}
    %     \begin{itemize}
    %         \item 某年学业先进个人
    %         \item 某年某奖学金某等奖
    %         \item 某大使
    %         \item 某年某奖学金某等奖
    %         \item 某年优秀团员称号
    %         \item 某年某称号
    %     \end{itemize}
    % \end{multicols}

    % % 其他
    % \section{\makebox[\widthof{\faInfo}][c]{\color{primary_color}{\faInfo}}\quad 其他}
    % \begin{itemize}
    %     \item 英语水平-CET6级xxx分
    %     \item 计算机几级证书
    %     \item xx几级证书
    %     \item 技术博客: 某网址
    %     \item 教师资格证：xxx
    %     \item 普通话证书：几级几等
    %     \item 文字排版：\LaTeX
    % \end{itemize}

\end{document}
